% =====================================================================
%  Harald Høffding, KORT OVERSIGT OVER DEN NYERE FILOSOFIS HISTORIE.
%  Fjerde, paany omarbejdede Udgave.  Kjøbenhavn og Kristiania: Gyldendalske Boghandel, Nordisk Forlag, 1910.
%  Forlagets Bogtrykkeri.  95 pp.
%
%  GENERATED; DO NOT EDIT BY HAND.  Built by ebook-method/kort-oversigt/ (rebuild.sh) from the scan, the
%  collation decisions and the patches there.
%
%  SOURCE: the National Library of Norway's scan (nb.no item 17246e2b574fbd978af6a80c866e4e4b,
%  URN:NBN:no-nb_digibok_2010102906103), ~/bibliotek/Høffding, Harald/1910-kort-oversigt-4udg.pdf (SCANS.tsv).
%  PDF page = printed page + 10.
%
%  EDITION (2026-10-17): the 1910 fourth edition, chosen by Hans; the only edition with an open scan (NB also holds
%  1906, 1919 and 1922, not open; Runeberg's kortfil is the same 1910 NB copy, unproofread OCR).  The SAGA /
%  Lindhardt og Ringhof e-book (ISBN 9788726363074, "Copyright 1898, 2020") is NOT the 1898 text but a LATER edition
%  (it gives Wundt's death, 1920; a second Tillæg „Vurderingsprincipets Historie“; its note cites Den menneskelige
%  Tanke 1910 where 1910's cites Mindre Arbejder; verb plurals modernised): a voting witness only where it agrees.
%
%  HOW ACCURACY WAS ESTABLISHED (2026-10-17; ebook-method/kort-oversigt/, README entry there):
%  (a) Draft: NB's own OCR (ALTO) laid over the page images (nb_scan.py), turned into LaTeX by rule (draft.py), with
%  running heads and sheet signatures dropped; headings from the Indhold; the one note placed at its mark.
%  (b) Word collation (collate/check.py, punctuation included) against an INDEPENDENT tesseract reading of the same
%  images: ~490 disagreements; 216 by the e-book's vote (rules.py), but where it voted for CAPITALS the word is in small
%  capitals in the print (SC, rules2.py: 87), 15 NB å for a, 5 umlauts; 19 by hand where three witnesses agree; 125 at
%  the image.  OPEN 0; italics OPEN 0.
%  (c) Blind crops (b1: two readers): the open checks, the e-book's non-word differences, italics, and an AUDIT of 30
%  rule decisions (30/30 right); controls 24/24.
%  (d) Structure read at the image (s1: 80 crops, two readers; s2: 14 crops, main session): half-title, title page, the
%  blank versos, Indhold, Fortale, the end leaves; every chapter and section head (and the short rule that closes each
%  part); the run-in names in small capitals (29 openings; 13 names mended); the α. β. γ. labels; doubtful paragraph
%  starts (8); the one footnote whole.
%  (e) Sweeps, read at line crops in the main session: quotation marks (the print uses „…“ only; NB dropped „ or read it
%  « or »: 45 places mended, every paragraph now balances), unknown words against the e-book and the other Høffding
%  transcriptions (13 mended: ü/ä/à/è/œ, l' read V), names half in small capitals (6 mended).  The whole text was also
%  compared with the (later-edition) e-book (ebookdiff.py: 1238 differences, mostly edition wording; its 32 non-word
%  suspects went to b1).
%  (f) A blind sampled page check (pagecheck.py): 6 pages from four strata (pp. 5, 12, 43, 44, 58, 82), 1,764 words:
%  0 word errors, 0 punctuation errors, 1 markup error (an italic title, mended); 95 % upper bound 1.7 word errors per
%  1000 words.
%  (g) The translators' 54 odd spots (tr/odd-*.tsv), read at line crops in the main session (o6-o8): 33 faults of the
%  transcription mended: the first line of p. 2, dropped as a running head (draft.py now keeps a line ending in a
%  hyphen); twelve doubled dashes; four specks read as stops or a comma; „done“ for donc; a doubled og; Göttingen,
%  Système (2); a lost paragraph break (p. 32); eight titles only partly italic.  Nine misprints of the print kept and
%  marked PRINTED AS IS.  The rest were questions of grammar, fact or period spelling, settled in translating.
%  NOT done: a whole-book reading at the image.
%
%  EDITORIAL: diplomatic.  Printed as is, and marked: „Er Sjælen er særegen“ (p. 1), „og og“ (p. 3), „Bevægelsee“ (p. 10),
%  „hviken“ (p. 29), „Skrifter“ without its stop (p. 32), „almægtigt Endeligt“ (p. 43), „Væsen Hvad“ (p. 44), „raade
%  Den“ (p. 48), „og og“ (p. 66), „Tænkeré“ (p. 68), „Aarhundres“ (p. 77), „Hypetese“ and „staar, En“ (p. 79),
%  „kvantitive“ and „concience“ (p. 83), „Livsvurdening“ (p. 94).  Kept as printed, unmarked:
%  Burckardt, Melanchton, Réné, Oldenbarnewelts, Darvin (Indhold), Darwin's death 1881, reduit, Holbach in italics.
%  The house dash --- for the print's —; the publisher's device on the title page is described, not drawn.
% =====================================================================
\documentclass[12pt,a4paper,openany]{book}
\usepackage[utf8]{inputenc}
\usepackage[T1]{fontenc}
\usepackage{libertinus}
\usepackage{textalpha}
\usepackage[danish]{babel}
\usepackage{setspace}
\usepackage[margin=1.2in]{geometry}
\usepackage{fancyhdr}
\setlength{\headheight}{28pt}
\addtolength{\topmargin}{-13.5pt}
\usepackage{hyperref}
\hypersetup{pdftitle={Kort Oversigt over den nyere Filosofis Historie (1910)}, pdfauthor={Harald Høffding}, hidelinks}
\pagestyle{fancy}
\fancyhf{}
\fancyhead[LE,RO]{\thepage}
\fancyhead[RE]{\textit{Harald Høffding: Kort Oversigt}}
\fancyhead[LO]{\textit{\leftmark}}
\fancypagestyle{plain}{\fancyhf{}\fancyhead[LE,RO]{\thepage}}
\setstretch{1.3}
\usepackage{marginnote}
\usepackage{graphicx}
\DeclareUnicodeCharacter{0254}{\reflectbox{c}}  % ɔ: the old Danish id-est mark (a reversed c)
\newcommand{\starmark}[1]{\ifcase#1\or*\or**\or***\fi)}   % the print marks its one note *)
\renewcommand{\thefootnote}{\starmark{\value{footnote}}}   % the print numbers its notes 1) 2) per page, a few *)
\newcommand{\opage}[1]{\marginnote{\footnotesize\textit{[#1]}}}
\newcommand{\apage}[1]{\marginnote{\footnotesize\textit{[#1]}}}
\newcommand{\whiteline}{\par\vspace{\baselineskip}\par}
% the printed Indhold: an essay or chapter line (title, page), a section line
\newcommand{\ientry}[2]{\par\smallskip\noindent #1\dotfill\,#2\par}
\newcommand{\isub}[2]{\par\noindent\hspace*{2em}#1\dotfill\,#2\par}
% the headings (no \addcontentsline: the print has its own Indhold, and a footnote in a heading breaks a toc entry)
\pdfstringdefDisableCommands{\renewcommand{\footnote}[2][]{}}
\newcommand{\prefacehead}[1]{\chapter*{#1}\markboth{Fortale}{Fortale}}
\newcommand{\sectionhead}[1]{\par\medskip\begin{center}\textbf{#1}\end{center}\smallskip}   % the print: centred, bold
\newcommand{\essayhead}[2][]{\clearpage\begin{center}{\Large #2}\end{center}\if\relax\detokenize{#1}\relax\markboth{#2}{#2}\else\markboth{#1}{#1}\fi\medskip}   % [running head] when the title carries a note
\newcommand{\essaysource}[1]{\begin{center}{\small #1}\end{center}\medskip}
\newcommand{\subhead}[1]{\par\medskip\begin{center}#1\end{center}\smallskip}
\newcommand{\introhead}[1]{\begin{center}{\large #1}\\[4pt]\rule{2em}{0.4pt}\end{center}\markboth{Indledning}{Indledning}\medskip}
\newcommand{\closerule}{\par\bigskip\begin{center}\rule{4em}{0.4pt}\end{center}\par}   % the short rule that closes each part of the print
\newcommand{\chapterheadone}[2]{\bigskip\begin{center}{\large #1}\\[4pt]\rule{2em}{0.4pt}\end{center}\markboth{#2}{#2}\medskip}   % VIII: numeral and title on one line
\newcommand{\chapterhead}[2]{\bigskip\begin{center}#1\\[3pt]{\large #2}\\[4pt]\rule{2em}{0.4pt}\end{center}\markboth{#2}{#2}\medskip}   % numeral, title in capitals, short rule
\newcommand{\romanhead}[1]{\par\medskip\begin{center}#1.\end{center}\smallskip}

\begin{document}

% ---------------------------------------------------------------------
%  FRONT MATTER (scan PDF 3-9), read at the image (structure pass s1, 2026-10-17).  PDF 4, 6, 8, 10 are blank;
%  nothing is printed after p. 95 (PDF 106-109: blank leaf, endpapers, board).
% ---------------------------------------------------------------------
\frontmatter
\thispagestyle{empty}
\vspace*{5cm}
\begin{center}
{\large KORT OVERSIGT OVER DEN\\[4pt] NYERE FILOSOFIS HISTORIE}
\end{center}
\clearpage
\begin{titlepage}
\centering
\vspace*{2cm}
{HARALD HØFFDING}\par
\vspace{6pt}\rule{2em}{0.4pt}\par
\vspace{1.2cm}
{\Huge KORT OVERSIGT OVER\\[6pt] DEN NYERE FILOSOFIS\\[6pt] HISTORIE}\par
\vspace{1cm}
{\small FJERDE, PAANY OMARBEJDEDE UDGAVE}\par
\vfill
{\small [the publisher's device: an oval wreath with the monogram GB]}\par
\vfill
{KJØBENHAVN OG KRISTIANIA}\par
{GYLDENDALSKE BOGHANDEL}\par
{NORDISK FORLAG}\par
\vspace{4pt}\rule{2em}{0.4pt}\par\vspace{4pt}
{\footnotesize FORLAGETS BOGTRYKKERI}\par
{\footnotesize 1910}\par
\end{titlepage}
\clearpage

% the Indhold (PDF 7): from NB's ALTO, read at the image (s1-03-01: every entry and page number as printed, Réné and Darvin included)
\chapter*{INDHOLD}
\markboth{Indhold}{Indhold}
\noindent\hfill{\small SIDE}\par
\ientry{Indledning. De filosofiske Problemer}{1}
\ientry{I. Den nyere Filosofis Tilbliven}{5}
\isub{1. Opdagelsen af det naturlige Menneske}{5}
\isub{2. Det nye Verdensbillede}{11}
\isub{3. Den nye Videnskab}{14}
\ientry{II. Systemdannelser}{17}
\isub{a. Réné Descartes}{18}
\isub{b. Thomas Hobbes}{21}
\isub{c. Baruch Spinoza}{26}
\isub{d. Gottfried Wilhelm Leibniz}{32}
\ientry{III. Erfaringsfilosofien}{35}
\isub{a. John Locke}{36}
\isub{b. George Berkeley}{38}
\isub{c. David Hume}{39}
\ientry{IV. Oplysningsfilosofien}{43}
\isub{1. Voltaire, Holbach og Rousseau}{44}
\isub{2. Wolff og Lessing}{48}
\ientry{V. Immanuel Kant og den kritiske Filosofi}{50}
\isub{Biografi og Karakteristik}{50}
\isub{1. Erkendelsesteori}{53}
\isub{2. Etik}{59}
\isub{3. Enhedslære paa æstetisk og biologisk Grundlag}{61}
\ientry{VI. Romantikens Filosofi}{63}
\isub{a. Fichte, Schelling og Hegel}{63}
\isub{b. Schopenhauer, Kierkegaard og Feuerbach}{67}
\ientry{VII. Positivismen}{74}
\isub{1. Saint Simon og Comte}{75}
\isub{2. Stuart Mill}{77}
\isub{3. Darvin og Spencer}{78}
\ientry{VIII. Nyere Problembehandlinger}{81}
\isub{a. Det psykologiske Problem}{82}
\isub{b. Tilværelsesproblemet}{85}
\isub{c. Erkendelsesproblemet}{88}
\isub{d. Vurderingsproblemet}{92}
\clearpage

% the Fortale til fjerde Udgave (PDF 9, unpaged): from NB's ALTO, read at the image (s1-04-01: one paragraph; the date italic, the signature bold)
\prefacehead{FORTALE TIL FJERDE UDGAVE}

I denne nye Udgave har jeg helt omarbejdet den lille Bog efter
de Erfaringer, jeg har gjort ved i en Række af Aar at benytte den
ved det propædeutiske Kursus. Jeg har udeladt Meget, som der
sjelden blev Brug for ved dette Kursus, men har til Gengæld optaget
biografiske Oplysninger, og jeg har desuden søgt at gøre
Fremstillingen mere afrundet og forstaaelig. Dette gælder ikke mindst
det sidste Afsnit, som skulde give en Forestilling om Filosofiens
nuværende Arbejde. --- Jeg bringer min bedste Tak til Hr. Dr.
Anton Thomsen for en Række værdifulde Bemærkninger, der ere
komne denne nye Udgave tilgode.

\textit{Den 28. Marts 1910.}

\hfill\textbf{Harald Høffding.}
\clearpage

\mainmatter
%%BODY%%

\end{document}
